\chapter{Wyniki eksperymentów}

\section{Wyniki strojenia hiperparametrów}

Wyniki strojenia hiperparametrów dla każdego z modeli implicit ALS zostały przedstawione w tabeli~\ref{tab:hiperparametry-als}. Tak jak wspomniano w metodologii parametr odpowiadający za ilość iteracji jest równy 30 w każdym z modeli.

\begin{table}[htbp]
\centering
\caption{Najlepsze hiperparametry wariantów ALS wybrane na zbiorze walidacyjnym}
\label{tab:hiperparametry-als}
\begin{tabular}{lrrrrr}
\toprule
Wariant & Czynniki & $\alpha$ & $\beta$ & Regularyzacja & \ndcg{} \\
\midrule
Liniowa & 80 & $9.948\cdot 10^{-7}$ & - & 0.003588 & 0.249\\
Logarytmiczna & 24 & 0.000571 & - & 0.130920 & 0.261\\
Logarytmiczna z saturacją & 24 & 0.024390 & - & 0.169714 & 0.261\\
Logarytmiczna z przesunięciem $\beta$ & 48 & 0.389185 & 108.284 & 0.010565 & 0.275\\
\bottomrule
\end{tabular}
\caption*{\footnotesize Źródło: opracowanie własne na podstawie wyników Optuny.}
\end{table}

Warto zwrócić szczególną uwagę na niskie wartości parametru
$\alpha$ odpowiadającego za skalowanie wagi pozytywnych interakcji. Jako standardową wartość tego hiperparametru zwykle przyjmuje się 1 \cite{sparkals,frederickson2022implicit}, a przeszukanie przeprowadza wsród większych wartości w celu zróżnicowania dlugich interakcji od ich braku. Konsekwencją tego jest bardzo małe wykorzystanie potencjalnych informacji o preferencjach zawartych w długości czasu gry. Taki wynik był spodziewany w przypadku liniowego przekształcenia ze względu na przypisywanie ogromnych wag pewności dla obserwacji z długiego ogona. Klasycznym rozwiązaniem tego problemu jest zastosowanie logarytmu, jednak pomimo poprawy \ndcg{} $\alpha$ nadal pozostała mała. Z tego powodu zdecydowano się na zastosowanie przekształcenia $c_{ui} = 1 + \alpha \frac{\log(1 + r_{ui})}{1 + \log(1 + r_{ui})}$ ograniczającego wartości zmiennej przed skalowaniem do przedziału $[0, 1)$. Nie miało to jednak znaczącego wpływu na metrykę \ndcg{} na zbiorze walidacyjnym, a dodatkowe eskperymenty pokazały marginalność różnic między tymi modelami. Wyniki poprzednich eksperymentów doprowadziły do sformułowania hipotezy, że niskie wartości \texttt{playtime\_forever} wprowadzały znaczny szum do modeli stosujących logarytmiczne przekształcenia, co w efekcie zmuszało dobór niewielkich $\alpha$ w procesie optymalizacji hiperparametrów. W celu zniwelowania tego zjawiska oraz nadania interakcjom z \texttt{playtime\_forever = 0} zastosowano $c_{ui} = 1 + \alpha \log(1 + r_{ui} + \beta)$. W procesie optymalizacji $\beta$ została ustalona na wartość 108,284, co umożliwiło skuteczne różnicowanie siły interakcji, bez ryzyka zdominowania procesu uczenia przez szum informacyjny krótkich sesji bądź przez dużą skalę wartości długich interakcji.
\section{Analiza funkcji wag pewności}
\section{Wyniki finalne wariantów ALS}
\section{Porównanie z modelami odniesienia}
\section{Trafność, pokrycie katalogu i obciążenie popularnością}
\section{Stabilność wyników, testy statystyczne i analiza wrażliwości}